\emph{Exponential observables \(\operatorname{Tr}\mathrm e^{\mathrm ik^\mu A_\mu}\) at \(N=2\).} Let \(O_k(A)=\operatorname{Tr}\mathrm e^{\mathrm ik^\mu A_\mu}\) with \(k\in\mathbb C^{10}\), at \(\gamma=0\), regulator (A), contours \(\mathcal R_u\) and \(\epsilon\to0^+\) last. For each fixed \(\epsilon>0\) and \(u<2/3\) the integral converges absolutely and is entire in \(k\), because on \(\mathcal R_u\) the growth of \(O_k\) is at most \(\mathrm e^{c(u,k)\|\tilde A\|}\) against a Gaussian regulator, and the fixed-\(\epsilon\) Stokes identity extends to it. At \(N=2\) the matrix \(k^\mu A_\mu\) has eigenvalues \(\pm\rho_X/2\) with \(\rho_X^2=k^\mu k^\nu\mathcal G_{\mu\nu}\), so \(O_k=2\cos(\rho_X/2)\) is branch-independent, and \(|2\cos(\omega/2)|^2=4(\cos^2(\operatorname{Re}\omega/2)+\sinh^2(\operatorname{Im}\omega/2))\). Along rank-one rays \(\tilde A_\mu=\rho_\mu H\), \(\operatorname{Tr}H^2=1\), the Euclidean weight \(\mathrm e^{-S_{\mathrm E}}\) is not exponentially damped and tends to a bounded factor along the corner (\(S_{\mathrm E}\) vanishes on the exact rank-one rays), the Pfaffian vanishes on the ray and is polynomial in the coordinates nearby, and the remaining factors are polynomial (numerical, in a reduced chart), while \(|O_k|\ge 2\sinh\!\bigl(b_u(\hat\rho;k)\rho/2\bigr)\) with \(b_u(\hat\rho;k)=|\operatorname{Im}(k^0\mathrm e^{-3\pi\mathrm iu/8}\hat\rho_0+k^i\mathrm e^{\pi\mathrm iu/8}\hat\rho_i)|\). At every \(u\in(0,1]\) the functional \(\hat\rho\mapsto\operatorname{Im}(k^0\mathrm e^{-3\pi\mathrm iu/8}\hat\rho_0+k^i\mathrm e^{\pi\mathrm iu/8}\hat\rho_i)\) is nonzero for every nonzero real \(k\) (at \(u=1\) its coefficients are \(-\sin(3\pi/8)k^0\) and \(\sin(\pi/8)\vec k\)), so the Euclidean side of the displayed limit identity diverges in modulus for every nonzero real \(k\); more generally it diverges for every \(k\) outside the slice on which \(k^0\mathrm e^{-3\pi\mathrm i/8}\) and all \(k^i\mathrm e^{\pi\mathrm i/8}\) are real, where \(|O_k|\le2\) and that side converges by the \(k<14\) Pfaffian window. For \(\operatorname{Im}k\ne0\) such that some commuting unit \(\hat\rho\) has \(|\operatorname{Im}(k\cdot\hat\rho)|>|\operatorname{Re}(k\cdot\hat\rho)|\), the commuting-ray contribution to the \(u=0\) integral has \(S_{\mathrm b}=0\) and grows along a sequence \(\epsilon\to0^+\) like \(\mathrm e^{c_2(k)/\epsilon}\) with \(c_2(k)>0\) (commuting-ray model; no transverse stationary-phase bound is claimed), and for real \(k\ne0\) the integrand at \(u=0\) carries \(|O_k|\le2\) pointwise, but no fixed-\(u\) domination in \(\epsilon\) is available at any \(u\in(0,2/3)\): the fixed-\(\epsilon\) modulus on \(\mathcal R_u\) grows without bound as \(\epsilon\to0\), numerically like \(\mathrm e^{c(k,u)/\sqrt\epsilon}\) on the sampled range. Hence within the class stated the displayed limit identity fails for every nonzero real \(k\) and for every non-real \(k\) outside the slice on which \(k^0\mathrm e^{-3\pi\mathrm i/8}\) and all \(k^i\mathrm e^{\pi\mathrm i/8}\) are real, holds at \(k=0\), and is unresolved on that slice at \(k\ne0\): there is no analyticity strip and no non-polynomial extension of the limit statement at \(\gamma=0\) off that slice, within the hypothesis class \(N=2\), regulator (A), contours \(\mathcal R_u\), \(\epsilon\to0^+\) last. The obstruction is the massless flat direction, and at fixed \(\gamma<0\) it is absent: the mass contributes on \(\mathcal R_u\) the real exponent \((\gamma/2g_{\mathrm{YM}}^2)[\sin(3\pi u/4)\operatorname{Tr}\tilde A_0^2+\sin(\pi u/4)\sum_{i=1}^{9}\operatorname{Tr}\tilde A_i^2]\), a Gaussian damping on every component that dominates \(\mathrm e^{c(k)\|\tilde A\|}\) uniformly for \(\epsilon\le\epsilon_0<|\gamma|/(2g_{\mathrm{YM}}^2)\), and the displayed limit identity then extends to \(O_k\) for every \(k\in\mathbb C^{10}\) (numerical). At fixed \(\gamma>0\) the \(u=1\) integral diverges in modulus even at \(k=0\) (numerical), so the contour family and the thimbles of~\cite{ChouNishimuraTripathi2025} act on different contour integrals and produce no competing value; the classicalization of~\cite{AsanoNishimuraPiensukYamamori2025} is tied there to a Lorentz-breaking infrared cutoff and bears neither on this paragraph nor on the statements above. Everything in this paragraph is for the hypothesis class stated in it, does not address \(N\to\infty\) taken first, and derives no signature, real-time response, background or continuation; Assumption~\ref{ass:lorentzian} is not discharged. Problem~E remains open and PROGRAMME-CONDITIONAL.
